\section{Conclusions}

OSR set out to answer a question no existing ranking answers. Where do people from my country actually go, and do they stay? Building and launching it taught me three things I did not expect.

First, shipping to a real audience is the strongest evaluation instrument I used, and the least structured. The launch comments carried sharper critique than any lab protocol would have elicited. One commenter stated the right-censoring problem correctly without knowing its name. But raw comments are not findings. They became design input only through coding, which also showed that most of the corpus is friend-tags and praise. Reach and signal arrive mixed. The work is in separating them.

Second, autobiographical design bought speed and cost perspective. Because I am the primary persona, requirements were never blocked on access to users, and the metric layer earned a depth a survey-driven process would not have justified. The blind spots sat exactly where my experience ran out. I underestimated geographic localisation. Commenters found administrative fragments listed as cities and an opaque name (``Staad'') that one commenter could only guess was a bus terminus. And I pitched explanations at my own statistical literacy, though the audience varies widely. Some readers summoned an AI bot to summarise the post; others critiqued denominator choices. Both corrections came from outside.

Third, the platform taught me a requirement I had never read anywhere. Views must survive as screenshots. Each takedown killed the circulating link. What persisted was the image. Screenshot-survivability (FR11) entered the requirements from watching distribution under censorship.

The process also has strengths: every requirement traces to a counted theme, benchmark finding, or PACT implication; every iteration-2 change is dated by a public commit; and the second evaluation fixed its decision rule before any data arrived.

The limits stated in Section~\ref{sec:pact} stand: one coder, one platform, and rates that describe only ORCID-registered, PhD-track scholars, while the comments show the largest interested groups, master's applicants and working professionals, sit outside that population. The sessions add their own caveats: three participants recruited from my own network, two of them in computer science; every session ran on desktop although the protocol called for one or two at phone width; and T3/T5 reached only the final session as my facilitation improved. Enough to expose the metric-semantics gap and the icon-labelling fix, and to hear all three name outcome data as the thing they could never find. Not enough to claim saturation.

Future work follows the requirements still open. The institution ranking needs a field-level breakdown before /ranking is promoted as a headline product. An R1 tier means little if it averages across disciplines. Every destination needs its sample size displayed, so Antarctica-style anomalies explain themselves. Side-by-side comparison would serve the offer decision the project started from. A Chinese-language option would stop the dominant audience reading the metrics through translation. Extending coverage beyond PhD-track researchers, the corpus's most-repeated request, needs data ORCID cannot provide, the longest road and the most necessary.

OSR ranks destinations by observed trajectories, conditioned on origin. No benchmarked product does that. It rests so far on one population, one dataset, and one platform's audience. The next iteration must widen all three.
